\section{模型的评价}

\subsection{模型的优点}

第一，三问共用逐年、逐季、逐地块的面积变量和土地约束体系。输出是可直接填写的完整面积表，并由独立程序逐项检查容量、适种性、重茬、豆类覆盖和水浇地模式，结果具有明确的复核路径。

第二，问题二和问题三将整数方案求解与随机评价分开。方案只求解一次，随后在固定的样本外情景中计算收益和风险，避免测试情景复制整数决策变量。正式运行因此能够在有限计算时间内完成方案生成与逐情景复核。

第三，主方法与基线使用同一批随机情景进行配对比较。平均收益、分位数和下尾均值的差异由同一概率场计算，减少了两组独立抽样造成的比较噪声。五个随机种子的结果还显示，测试范围内主方案的平均收益方向保持一致。

第四，问题三同时使用收益偏差和 41 维作物累计面积相似度判断结构变化。弱、中、强三档的宏观面积相似度均高于 80\%，因此将较低的逐地块重合度解释为微观调度差异，而不是直接判为市场结构跳变。

\subsection{模型的不足}

第一，问题一的销量由 2023 年产量构造，价格取区间中点。该设定适合形成确定性比较基准，但不能替代多年市场预测，滞销模式的未闭合 Gap 也意味着对应收益不能称为精确全局最优值。

第二，问题二和问题三的概率分布由题目区间和建模设定给出，而不是由多年历史样本估计。200 个情景中没有出现亏损，只能说明在已测试的模拟范围内未观察到亏损，不能解释为现实经营中的亏损概率为零。

第三，问题三的替代、互补和经济变量相关结构均为透明的模拟假设。相关矩阵核验只能验证随机生成和映射的数值实现，不能据此识别真实交叉价格弹性或因果关系。

第四，模型没有纳入极端天气、病虫害、土地质量变化、库存和跨季销售，也没有刻画劳动力峰值与地块间运输成本。因此，相关结论的适用范围限于题目给定的土地条件、参数区间和销售规则。
